\chapter{Fazit}
\label{kap:8}

Dieses abschließende Kapitel fasst die Ergebnisse der vorliegenden Arbeit zusammen (Abschnitt~\ref{sec:fazit_zusammenfassung}), reflektiert deren Aussagekraft kritisch im Hinblick auf methodische und externe Einschränkungen (Abschnitt~\ref{sec:fazit_reflexion}) und skizziert Ansatzpunkte für künftige Weiterentwicklungen (Abschnitt~\ref{sec:fazit_ausblick}).

% ─────────────────────────────────────────────────────────────────────────────
\section{Zusammenfassung der Ergebnisse}
\label{sec:fazit_zusammenfassung}

\subsection{Beitrag zur übergeordneten Forschungsfrage}
\label{sec:fazit_ff}

Mit dem \ac{ACE}-System wurde ein \ac{DSR}-Artefakt entwickelt, das regelbasierte Prüfung, Interaktionsanalyse, Quellcodekontext und \ac{LLM}-gestützte Auswertung in einer Pipeline zusammenführt. Die Hauptevaluation erfolgte auf 90 Benchmarkläufen (3 Modelle $\times$ 10 Runs je Testsuite) über die drei synthetischen Testsuiten test-12, test-50 und test-100. Dabei dient test-12 vor allem der Kalibrierung, test-50 als Zwischenstufe und test-100 als Belastungsprüfung.

Im Rahmen dieser synthetischen Evaluation kann \textbf{TF1} bestätigt werden. Die kombinierte Pipeline erzielt auf der test-100-Suite mit \texttt{qwen3:32b} einen konsistenten Recall von 63\,/\,100~Violations und liegt damit deutlich über der Tool-only-Baseline von 34\,/\,100. Damit zeigt sich, dass die Ergänzung klassischer Tools um \ac{LLM}-gestützte Analyse zusätzliche, insbesondere semantische und kontextabhängige Barrieren erschließen kann. Zugleich macht das Ergebnis deutlich, dass auch das beste untersuchte Modell weit von vollständiger Abdeckung entfernt bleibt. Die Arbeit zeigt daher keinen vollautomatischen Ersatz manueller Accessibility-Prüfung, sondern eine substanzielle Erweiterung bestehender Toolketten.

\textbf{TF2} wird durch die Ergebnisse überwiegend gestützt. Auf der test-12-Suite erreicht \texttt{qwen3:32b} für 11\,/\,12~Violations eine Empfehlungsqualität der Stufe~2, also konkrete, entwicklungsnahe Fix-Vorschläge mit Dateibezug und umsetzbarem Codekontext. Dies spricht dafür, dass die Zusammenführung von Tool-Findings, Interaktionsdaten und Quellcodekontext die Qualität der Handlungsempfehlungen deutlich verbessert. Die Aussage ist jedoch bewusst vorsichtig zu interpretieren: Da keine isolierte Ablationsstudie ohne Quellcodekontext durchgeführt wurde, lässt sich der Effekt nicht kausal einem einzelnen Kontextbaustein zuschreiben, sondern nur der Gesamtarchitektur der Pipeline. Detailwerte finden sich in den Tabellen~\ref{tab:recall_matrix}, \ref{tab:modell_summary_12}, \ref{tab:recall_matrix_100} und \ref{tab:empfehlungsqualitaet}.

\subsection{Erfüllung der funktionalen und nicht-funktionalen Anforderungen}
\label{sec:fazit_anforderungen}

Die fünf funktionalen Anforderungen \ac{FA}-01 bis \ac{FA}-05 wurden im Sinne des entwickelten \ac{PoC} vollständig umgesetzt. Axe-core übernimmt die regelbasierte Basiserkennung (\ac{FA}-01), Playwright erfasst zustandsabhängige bzw. interaktive Barrieren (\ac{FA}-02), grep-gestützte Quellcodeanalyse erweitert die Sicht auf statische Muster (\ac{FA}-03), und die zweistufige \ac{LLM}-Verarbeitung erzeugt priorisierte, kontextsensitive Handlungsempfehlungen (\ac{FA}-04). Mit dem strukturierten \acs{JSON}-Report liegt zudem ein maschinenlesbares Ausgabeformat vor, das sich für Nachverarbeitung und potenzielle Integration in weitere Werkzeuge eignet (\ac{FA}-05).

Differenzierter fällt die Bewertung der nicht-funktionalen Anforderungen aus. \ac{NFA}-01 (Lokal-First / Datenschutz) wird durch die Ollama-basierte Ausführung vollständig erfüllt: Weder Quellcode noch \ac{DOM}-Inhalte verlassen das lokale System. \ac{NFA}-02 (Laufzeit $<$ 10~Min.) wird hingegen nur teilweise erreicht. Unter Kalibrierungsbedingungen der test-12-Suite erfüllen \texttt{qwen2.5-coder:7b} und \texttt{qwen2.5-coder:14b} die Anforderung; für \texttt{qwen3:32b} ist sie dort nur eingeschränkt erfüllt. Auf der test-100-Suite überschreiten hingegen alle drei Modelle die Zehn-Minuten-Grenze deutlich (Ø~765--2\,616\,s). Daraus ergibt sich eine wichtige Einsatzkontext-Erkenntnis: In der vorliegenden Form eignet sich \ac{ACE} bei größeren Violation-Mengen eher für Release-nahe Prüfungen, Nightly-Analysen oder gezielte Qualitätssicherungsphasen als für eng getaktete per-Commit-Läufe. Eine Nutzung im Sprint- oder Komponenten-Kontext erscheint zwar plausibel, wurde in dieser Arbeit jedoch nicht separat empirisch untersucht und bleibt daher eine begründete Extrapolation, keine direkte Messung. Die entsprechende Kontextdifferenzierung wird in Tabelle~\ref{tab:einsatzempfehlungen} sowie in den Anhang-Tabellen~\ref{tab:suite_vergleich} und~\ref{tab:modell_summary_100} aufgegriffen.

\ac{NFA}-03 (Wartbarkeit) wird durch die modulare Trennung von Detection, Context-Building und Synthesis unterstützt. \ac{NFA}-04 (Reproduzierbarkeit) wird teilweise erreicht: Temperature~0,05 und schemaorientiertes Prompting fördern stabile Ausgaben, die Varianz steigt jedoch mit der Suite-Größe und ist modellabhängig. \texttt{qwen2.5-coder:14b} zeigt auf test-12 mit $\sigma = 0{,}42$ die höchste Konsistenz, während auf test-100 auch das stärkste Modell keine vollständig stabile Ausgabe erreicht. Die Reproduzierbarkeit ist damit für einen experimentellen \ac{PoC} ausreichend, aber nicht in dem Sinne gegeben, dass identische qualitative Resultate über alle Run-Kontexte hinweg garantiert wären. Eine tabellarische Gesamtübersicht findet sich in Tabelle~\ref{tab:fa_nfa_erfuellung} im Anhang.

% ─────────────────────────────────────────────────────────────────────────────
\section{Kritische Reflexion}
\label{sec:fazit_reflexion}

\subsection{Grenzen der externen Validität}
\label{sec:fazit_extern}

Die belastbarste Evidenz dieser Arbeit stammt aus den drei synthetischen Testsuiten. Diese wurden gezielt so konstruiert, dass unterschiedliche Verstoßkategorien, verschiedene Schwierigkeitsgrade und Skalierungsstufen abgebildet werden. Gleichwohl bleiben synthetisch eingebrachte Violations in ihrer Struktur kontrollierter und oft klarer isolierbar als Barrieren in realen Anwendungssystemen, in denen technische Altlasten, Projektkonventionen, Framework-Effekte und historisch gewachsene UI-Muster zusammenwirken. Der auf test-100 mit \texttt{qwen3:32b} erreichte Recall von 63\,\% ist daher als belastbarer Laborwert innerhalb der gewählten Evaluationslogik zu verstehen, nicht als direkt übertragbarer Erwartungswert für produktive Systeme.

Die ergänzende Erprobung am \ac{BWEC} erweitert die Aussagekraft der Arbeit, ersetzt diese Einschränkung jedoch nicht. Sie zeigt, dass das System grundsätzlich auf ein reales Untersuchungsobjekt übertragen werden kann und dort plausible, datei- und zeilennahe Befunde erzeugt. Zugleich macht gerade der \ac{BWEC}-Block deutlich, dass die Skalierungsprobleme der Priorisierungsphase in realen Codebasen stärker hervortreten als in den synthetischen Suiten: Bei \texttt{maxfiles}~=~500 treten Truncation, kollabierte Outputs und fehlgeschlagene Priorisierung deutlich sichtbar auf. Der \ac{BWEC} ist damit vor allem als Feldtest der praktischen Einsetzbarkeit und Systemgrenzen zu verstehen, nicht als quantitative Validierung des Recall.

Hinzu kommt, dass die Laufzeitmessungen ausschließlich auf der verwendeten Hardwarekonfiguration und ohne \ac{GPU}-Beschleunigung erhoben wurden. Es ist plausibel, dass leistungsfähigere Inferenzumgebungen die Laufzeiten reduzieren würden; dies wurde jedoch in der vorliegenden Arbeit nicht experimentell untersucht. Aussagen zur Erfüllbarkeit von \ac{NFA}-02 auf alternativer Hardware bleiben daher hypothetisch.

\subsection{Methodische Einschränkungen}
\label{sec:fazit_intern}

Methodisch wird die Aussagekraft der Ergebnisse vor allem durch drei Faktoren begrenzt: Token-Budget, Operationalisierung der Erkennungsmetrik und begrenzte Tiefenvalidierung der generierten Reports.

Erstens beeinflusst das verfügbare Token-Budget die Priorisierungsqualität systematisch und modellabhängig. Auf test-100 schöpft \texttt{qwen2.5-coder:14b} in allen 10~Runs das Output-Limit von 12\,288~Tokens vollständig aus; \texttt{qwen2.5-coder:7b} in 80\,\% der Runs. Eine Erhöhung von \texttt{num\_predict} könnte diese Abschneideeffekte zwar reduzieren, würde jedoch zugleich die Laufzeit erhöhen und die Zielerreichung von \ac{NFA}-02 weiter verschlechtern. \texttt{qwen3:32b} liefert auf test-100 die methodisch robusteste Basis, weil nur 1\,/\,10~Runs trunkiert wurde; auch dieses Modell löst das Skalierungsproblem jedoch nicht vollständig.

Zweitens beruht die Erkennungsbewertung auf einer Zuordnung zwischen Ground-Truth-Violations und im finalen Report präsenten Findings. Diese Operationalisierung ist für die Arbeit praktikabel, bleibt aber konservativ: Semantisch äquivalente, anders formulierte oder gebündelte Findings können zu Unterschätzungen führen. Hinzu kommt, dass der Recall über eine Konsistenzschwelle ($>\,5/10$~Runs) operationalisiert wird. Diese Schwelle ist für die Untersuchung sinnvoll, markiert aber eine Modellierungsentscheidung und keine naturgegebene Grenze. Die Ergebnisse sind deshalb deskriptiv belastbar, aber nicht im Sinne inferenzstatistischer Tests verallgemeinerbar.

Drittens wurde keine vollständige Precision-Evaluation über alle 30~Runs der großen Suite durchgeführt. Die in den Überdetektions-Tabellen ausgewiesenen Quoten messen ausschließlich das Verhältnis von \ac{LLM}-Detektor-Rohoutput zur Ground-Truth-Größe und erlauben keine direkte Aussage über die inhaltliche Richtigkeit jedes einzelnen Report-Eintrags. Die ergänzende Stichprobenprüfung für \texttt{qwen3:32b} auf test-100 zeigt allerdings, dass echte Halluzinationen selten auftreten: In Run~01 wurden keine bestätigten Halluzinationen identifiziert; in Run~03 wurde eine Halluzination festgestellt. Der größere Teil der Überdetektion entsteht somit nicht durch frei erfundene Befunde, sondern durch Quell-Duplikate, Mehrfachinstanzen und nicht eindeutig zuordenbare Paraphrasen. Für einen produktiven Einsatz bedeutet dies dennoch, dass der Report nicht ungeprüft als abschließendes Audit-Ergebnis verwendet werden sollte, sondern eine nachgelagerte Plausibilisierung bzw. Deduplikation erforderlich bleibt.

Aus demselben Grund ist die vertiefende Fehlerursachenanalyse auf test-100 mit \texttt{qwen3:32b} konzentriert. Diese Modellkonfiguration weist innerhalb der großen Suite die geringste Truncation auf und stellt damit die methodisch sauberste Grundlage für die qualitative Nachanalyse dar. Auch diese Entscheidung reduziert Verzerrungen, beseitigt sie jedoch nicht vollständig.

% ─────────────────────────────────────────────────────────────────────────────
\section{Ausblick}
\label{sec:fazit_ausblick}

\subsection{Kurzfristige Erweiterungen}
\label{sec:ausblick_kurz}

Kurzfristig erscheinen vier Weiterentwicklungen besonders naheliegend. Erstens sollte der \ac{BWEC}-Block um eine systematische manuelle Plausibilisierung ergänzt werden. Damit ließe sich der reale Feldtest methodisch deutlich stärken, etwa durch einen Expertenabgleich der generierten Findings mit dem tatsächlichen Quellcode sowie durch eine zusätzliche referenzielle Baseline ohne \ac{LLM}-Anteile.

Zweitens sollte die Playwright-Komponente um gezielte Interaktionsprüfungen erweitert werden, insbesondere für Fokus-Traps, Dialog-Navigation und komplexere Keyboard-Flows. Gerade die test-12- und test-100-Ergebnisse zeigen, dass einzelne interaktive Barrieren derzeit noch nicht zuverlässig regel- oder zustandsbasiert abgefangen werden und teilweise erst in der \ac{LLM}-Stufe sichtbar werden.

Drittens wäre eine nachgelagerte Deduplikations- und Konsolidierungslogik sinnvoll. Ein Teil der Überdetektion entsteht nicht durch falsche Inhalte, sondern durch mehrfach berichtete oder unterschiedlich formulierte Varianten desselben Problems. Eine regelbasierte Konsolidierung zwischen axe-core-, Playwright-, grep- und \ac{LLM}-Findings könnte die Report-Qualität erhöhen und zugleich das Prompt-Budget der Priorisierungsphase entlasten.

Viertens sollte die manuelle Stichprobenprüfung als formalisierter Bestandteil des Workflows etabliert werden. Eine strukturierte Prüfung eines definierten Anteils der Report-Einträge -- etwa der nur in einzelnen Runs auftretenden oder ausschließlich \ac{LLM}-basierten Findings -- würde es erlauben, Halluzinationsraten, Deduplikationsbedarf und praktische Nutzbarkeit systematischer zu quantifizieren.

\subsection{Mittelfristige Weiterentwicklungen}
\label{sec:ausblick_mittel}

Mittelfristig stehen vor allem Verbesserungen der Skalierbarkeit im Vordergrund. Erstens könnte die Priorisierungsphase hierarchisch umgebaut werden: Anstatt alle Findings in einen einzigen großen Prompt zu überführen, wären mehrstufige Zusammenfassungs- oder Clusterverfahren denkbar, bei denen zunächst teilgruppenspezifische Priorisierungen und erst anschließend eine Gesamtkonsolidierung erfolgen. Gerade die \ac{BWEC}-Messungen legen nahe, dass hier das zentrale technische Bottleneck der aktuellen Architektur liegt.

Zweitens könnte die Zustandsprüfung stärker ausgebaut werden, indem axe-core nicht nur auf einem Endzustand, sondern auf mehreren systematisch durchlaufenen \ac{DOM}-Zuständen ausgeführt wird. Dies würde die Abdeckung dynamischer Barrieren erhöhen, zugleich aber die Laufzeit und die Anzahl zu priorisierender Findings weiter steigern. Eine solche Erweiterung müsste daher mit intelligenten Sampling- oder Priorisierungsstrategien kombiniert werden.

Drittens wäre die Evaluation größerer lokal ausführbarer Modelle auf geeigneterer Hardware ein naheliegender nächster Schritt. Die vorliegenden Ergebnisse deuten darauf hin, dass Modellgröße und Modellqualität einen erheblichen Einfluss auf Recall, Priorisierung und Fix-Konkretheit haben. Ob dieser Zusammenhang auch jenseits der hier untersuchten Modellklassen linear fortbesteht, bleibt jedoch offen und müsste empirisch geprüft werden.

Viertens bietet der strukturierte \acs{JSON}-Report eine gute Grundlage für nachgelagerte Automatisierung. Denkbar wäre etwa eine Toolkette, in der priorisierte Findings nicht nur dokumentiert, sondern direkt in Issue-Tracker überführt oder als Kontext für nachgelagerte Code-Assistenzsysteme genutzt werden.

\subsection{Langfristige Forschungsperspektiven}
\label{sec:ausblick_lang}

Langfristig erscheint insbesondere eine stärkere Spezialisierung des Modells vielversprechend. Ein domänenspezifisch angepasstes Modell, das mit Accessibility-relevanten React-/TypeScript-Beispielen und typischen \ac{BITV}-Kontexten trainiert oder feinabgestimmt wurde, könnte systematische Fehlklassifikationen reduzieren und zugleich die Priorisierungslogik stabilisieren. Damit würde sich der Fokus von einer reinen Modellgrößenstrategie hin zu einer stärker wissens- und domänenspezifischen Optimierung verschieben.

Darüber hinaus ist eine graduelle Integration in Entwicklungsprozesse denkbar. Die vorliegende Arbeit zeigt allerdings ebenso klar, dass eine vollständige Analyse großer Codebasen in jeder Pipeline-Ausführung derzeit nicht realistisch ist. Realistischer wäre daher ein gestuftes Betriebsmodell: leichte Prüfungen auf Komponenten- oder Änderungsbasis während der Entwicklung, tiefere Analysen in qualitätssichernden Release- oder Nightly-Kontexten. Gerade in dieser differenzierten Form könnte \ac{ACE} zu einem praktikablen Baustein kontinuierlicher Barrierefreiheitsprüfung werden.

In einer weiterentwickelten Form ließe sich daraus ein \textit{Continuous-Accessibility-Loop} ableiten: automatische Erkennung, konsolidierte Priorisierung, unterstützte Behebung und erneute Prüfung. Der Beitrag eines solchen Systems läge dann nicht in der vollständigen Automatisierung von Accessibility-Audits, sondern in der nachhaltigen Reduktion manuellen Prüfaufwands und in der frühzeitigen Sichtbarmachung relevanter Barrieren im Entwicklungsprozess.